% ============================================================
\chapter{Central Limit Theorem}
\label{ch:clt}
% ============================================================

\begin{intuition}
The central limit theorem (CLT) is arguably the most celebrated result in
probability theory. It states that the normalised sum of i.i.d.\ random variables
converges in distribution to the normal law, regardless of the common distribution,
provided the variance is finite.
\end{intuition}

% ------------------------------------------------------------
\section{Statement of the Theorem}
% ------------------------------------------------------------

\begin{theorem}[Central limit theorem]
\label{thm:clt}
Let $(X_n)_{n \geq 1}$ be i.i.d.\ with $\E[X_1] = \mu$ and
$\Var(X_1) = \sigma^2 \in (0, \infty)$. Set $S_n = \sum_{i=1}^n X_i$. Then
\[
    Z_n = \frac{S_n - n\mu}{\sigma\sqrt{n}}
    = \frac{\bar{X}_n - \mu}{\sigma/\sqrt{n}}
    \xrightarrow{\mathcal{L}} \mathcal{N}(0, 1).
\]
That is, for every $x \in \R$:
\[
    \lim_{n \to \infty}\PP(Z_n \leq x) = \Phi(x)
    = \frac{1}{\sqrt{2\pi}}\int_{-\infty}^{x} e^{-t^2/2}\,dt.
\]
\end{theorem}

\begin{intuition}[title=\textbf{Why is the normal distribution so ubiquitous?}]
The CLT explains why so many natural phenomena are approximately normal: any
quantity resulting from the sum of many small independent contributions will be
approximately Gaussian.
\end{intuition}

% ------------------------------------------------------------
\section{Proof via Characteristic Functions}
% ------------------------------------------------------------

We refer to Chapter~11 for the full theory of characteristic functions. Here we
sketch the proof, taking the L\'evy continuity theorem for granted.

\begin{proof}[Proof of the CLT]
WLOG set $\mu = 0$ (otherwise replace $X_i$ by $X_i - \mu$). Let
$\varphi(t) = \E[e^{itX_1}]$ be the common characteristic function.

By assumption, $\varphi(0) = 1$, $\varphi'(0) = i\E[X_1] = 0$, and
$\varphi''(0) = -\E[X_1^2] = -\sigma^2$. Taylor expansion gives:
\[
    \varphi(t) = 1 - \frac{\sigma^2 t^2}{2} + o(t^2) \quad (t \to 0).
\]

The characteristic function of $Z_n = S_n/(\sigma\sqrt{n})$ is:
\begin{align*}
    \varphi_{Z_n}(t) &= \E\!\left[\exp\!\left(it\cdot\frac{S_n}{\sigma\sqrt{n}}\right)\right]
    = \left[\varphi\!\left(\frac{t}{\sigma\sqrt{n}}\right)\right]^n \\
    &= \left[1 - \frac{\sigma^2}{2}\cdot\frac{t^2}{\sigma^2 n}
       + o\!\left(\frac{1}{n}\right)\right]^n
    = \left[1 - \frac{t^2}{2n} + o\!\left(\frac{1}{n}\right)\right]^n.
\end{align*}
Using $\lim_{n\to\infty}(1 + a/n)^n = e^a$:
\[
    \varphi_{Z_n}(t) \xrightarrow[n\to\infty]{} e^{-t^2/2}.
\]
Since $e^{-t^2/2}$ is the characteristic function of $\mathcal{N}(0,1)$,
L\'evy's continuity theorem gives
$Z_n \xrightarrow{\mathcal{L}} \mathcal{N}(0,1)$.
\end{proof}

% ------------------------------------------------------------
\section{Practical Applications}
% ------------------------------------------------------------

\subsection{Normal approximation to the binomial}

\begin{corollary}[De Moivre--Laplace theorem]
If $S_n \sim \mathrm{Bin}(n, p)$ with $0 < p < 1$, then
\[
    \frac{S_n - np}{\sqrt{np(1-p)}} \xrightarrow{\mathcal{L}} \mathcal{N}(0,1).
\]
\end{corollary}

\begin{example}
A fair coin is tossed 100 times. What is the probability of getting between 45
and 55 heads?

With $S_{100} \sim \mathrm{Bin}(100, 0.5)$, $\mu = 50$, $\sigma = 5$:
\begin{align*}
    \PP(45 \leq S_{100} \leq 55)
    &= \PP\!\left(\frac{45-50}{5} \leq Z \leq \frac{55-50}{5}\right)
    \approx \PP(-1 \leq Z \leq 1) \\
    &= 2\Phi(1) - 1 \approx 0.6827.
\end{align*}
\end{example}

\begin{warning}[title=\textbf{Continuity correction}]
To improve the approximation, use the \textbf{continuity correction}:
\[
    \PP(a \leq S_n \leq b) \approx
    \Phi\!\left(\frac{b+0.5-np}{\sqrt{np(1-p)}}\right)
    - \Phi\!\left(\frac{a-0.5-np}{\sqrt{np(1-p)}}\right).
\]
\end{warning}

\subsection{Confidence intervals}

\begin{corollary}[Confidence interval for the mean]
If $(X_i)$ are i.i.d.\ with $\E[X_1] = \mu$ and known $\Var(X_1) = \sigma^2$,
a \textbf{95\% confidence interval} for $\mu$ is
\[
    \left[\bar{X}_n - 1.96\frac{\sigma}{\sqrt{n}},\;
    \bar{X}_n + 1.96\frac{\sigma}{\sqrt{n}}\right].
\]
\end{corollary}

\begin{keyformulas}[title=\textbf{Common quantiles for confidence intervals}]
\begin{center}
\renewcommand{\arraystretch}{1.3}
\begin{tabular}{ccc}
    \toprule
    Confidence level & $\alpha$ & $z_{\alpha/2}$ \\
    \midrule
    90\% & 0.10 & 1.645 \\
    95\% & 0.05 & 1.960 \\
    99\% & 0.01 & 2.576 \\
    \bottomrule
\end{tabular}
\end{center}
\end{keyformulas}

\subsection{Normal approximation to the Poisson}

\begin{corollary}
If $X \sim \mathrm{Poisson}(\lambda)$ with $\lambda$ large, then
$(X - \lambda)/\sqrt{\lambda} \xrightarrow{\mathcal{L}} \mathcal{N}(0,1)$.
\end{corollary}

% ------------------------------------------------------------
\section{Multivariate CLT}
% ------------------------------------------------------------

\begin{theorem}[Multivariate CLT]
Let $(X_n)$ be i.i.d.\ random vectors in $\R^d$ with $\E[X_1] = \boldsymbol{\mu}$
and $\Cov(X_1) = \Sigma$. Then
\[
    \sqrt{n}\,(\bar{X}_n - \boldsymbol{\mu})
    \xrightarrow{\mathcal{L}} \mathcal{N}_d(0, \Sigma).
\]
\end{theorem}

% ------------------------------------------------------------
\section{The Delta Method}
% ------------------------------------------------------------

\begin{theorem}[Delta method]
\label{thm:delta}
If $\sqrt{n}(\bar{X}_n - \mu) \xrightarrow{\mathcal{L}} \mathcal{N}(0, \sigma^2)$
and $g$ is differentiable at $\mu$ with $g'(\mu) \neq 0$, then
\[
    \sqrt{n}(g(\bar{X}_n) - g(\mu))
    \xrightarrow{\mathcal{L}} \mathcal{N}(0, \sigma^2[g'(\mu)]^2).
\]
\end{theorem}

\begin{example}
If $\bar{X}_n \to p$ and $g(x) = x(1-x)$, then $g'(p) = 1-2p$ and
\[
    \sqrt{n}(\bar{X}_n(1-\bar{X}_n) - p(1-p))
    \xrightarrow{\mathcal{L}} \mathcal{N}(0,\, p(1-p)(1-2p)^2).
\]
\end{example}

% ------------------------------------------------------------
\section{Rate of Convergence: Berry--Esseen}
% ------------------------------------------------------------

\begin{theorem}[Berry--Esseen]
Under the CLT hypotheses, if $\E[\abs{X_1}^3] < \infty$:
\[
    \sup_{x \in \R}\abs{\PP(Z_n \leq x) - \Phi(x)}
    \leq \frac{C\,\E[\abs{X_1-\mu}^3]}{\sigma^3\sqrt{n}},
\]
where $C$ is a universal constant ($C \leq 0.4748$).
\end{theorem}

\begin{remark}
Berry--Esseen shows that convergence in the CLT is of order $1/\sqrt{n}$.
This justifies using the normal approximation once $n$ is moderately large
(typically $n \geq 30$).
\end{remark}

% ------------------------------------------------------------
\section{Exercises}
% ------------------------------------------------------------

\begin{exercise}
Let $(X_i)$ be i.i.d.\ with $\E[X_1] = 2$, $\Var(X_1) = 9$. Approximate
$\PP(S_{100} > 220)$.
\end{exercise}

\begin{exercise}
A fair die is rolled 360 times. Approximate the probability that the sum exceeds
1300.
\end{exercise}

\begin{exercise}
Let $S_n \sim \mathrm{Bin}(400, 0.4)$. Compute approximately
$\PP(S_n \geq 180)$ with and without continuity correction.
\end{exercise}

\begin{exercise}
A physical quantity is measured $n = 50$ times. The measurements are i.i.d.\ with
variance $\sigma^2 = 4$. What is the 99\% confidence interval for the mean?
\end{exercise}

\begin{exercise}
How many polls are needed to estimate a proportion $p$ to within 2\% with 95\%
confidence? (The value of $p$ is unknown.)
\end{exercise}

\begin{exercise}
Using the delta method, find the limiting distribution of
$\sqrt{n}(\ln\bar{X}_n - \ln\mu)$ when the $X_i$ are i.i.d.\ positive with
$\E[X_1] = \mu > 0$ and $\Var(X_1) = \sigma^2$.
\end{exercise}

\begin{exercise}
Verify the CLT for $X_i \sim \mathrm{Exp}(\lambda)$ by directly computing the
characteristic function of $Z_n$ and showing it converges to $e^{-t^2/2}$.
\end{exercise}

\begin{exercise}
Let $(X_i)$ be i.i.d.\ $\mathrm{Poisson}(4)$. Approximate
$\PP\bigl(\sum_{i=1}^{100} X_i \leq 380\bigr)$.
\end{exercise}
